\chapter{Nerve Block}

\section*{What Is This Test or Procedure?}
A nerve block is an injection of local anesthetic, often with corticosteroid, near a nerve or nerve plexus to relieve pain, diagnose its source, or provide anesthesia.

\section*{Alternative Names}
Nerve block injection; peripheral nerve block; sympathetic block; plexus block

\section*{Principle}
Local anesthetic placed adjacent to a nerve temporarily interrupts nerve conduction, relieving pain and providing anesthesia. Corticosteroid reduces local inflammation. Fluoroscopy or ultrasound guidance targets the nerve precisely.

\section*{History}
Nerve blocks evolved from early regional anesthesia in the late 19th century into a mainstay of chronic pain management and surgical anesthesia in the 20th century.

\section*{Why This Test or Procedure Is Ordered}
Performed for chronic pain conditions (e.g., facet arthropathy, trigeminal neuralgia, complex regional pain syndrome), to provide surgical or obstetric anesthesia, and as a diagnostic tool to identify pain generators.

\section*{Is This a Primary or Secondary Test or Procedure}
Primary interventional treatment for selected pain syndromes; secondary to conservative management.

\section*{How This Test or Procedure Is Performed}
Performed in an outpatient or surgical setting. The target is identified by palpation, fluoroscopy, or ultrasound, the skin is prepared and anesthetized, and the needle is advanced to the nerve. Contrast and test doses may confirm placement before anesthetic injection.

\section*{What Preparation Is Needed}
Anticoagulants are usually held, informed consent is obtained, and patients are monitored during and briefly after injection.

\section*{Contraindications and Precautions}
Coagulopathy, anticoagulation, active infection, and allergy to local anesthetic are contraindications. Precaution: needle placement near major vessels, nerves, and the neuraxis requires image guidance.

\section*{Risks}
Transient numbness or weakness, bleeding, infection, nerve injury, systemic local anesthetic toxicity, and vasovagal reactions.

\section*{What Happens After the Test or Procedure}
The block resolves over hours; corticosteroid effects may take days. Patients are observed and warned about temporary weakness if a motor block occurred.

\section*{Understanding Your Results}
Temporary relief confirms the nerve as the pain source; prolonged relief indicates a therapeutic response.

\section*{Normal Results}
Significant pain relief lasting weeks to months without complication.

\section*{Follow-up Tests and Procedures}
Serial assessments; repeated blocks or referral to surgery may be considered based on response.

\section*{References}
\begin{enumerate}
  \item MedlinePlus. Nerve block. U.S. National Library of Medicine; 2025.
  \item Current Medical Diagnosis and Treatment. McGraw-Hill; 2025.
\end{enumerate}

\clearpage
